\section{Transverse Hensel specialization in the Pell--Lucas packet}
\label{sec:pell-polynomial-hensel}

Let
\[
 F_0(T)=0,\quad F_1(T)=1,\quad
 F_{n+2}(T)=T F_{n+1}(T)+F_n(T)
\]
and
\[
 L_0(T)=2,\quad L_1(T)=T,\quad
 L_{n+2}(T)=T L_{n+1}(T)+L_n(T).
\]
For $(1+\sqrt2)^n=A_n+B_n\sqrt2$ one has
$F_n(2)=B_n$ and $L_n(2)=2A_n$.  Recent work on powerful
Fibonacci and Horadam polynomials clarifies polynomial multiplicities
\cite{BatesEtAl2026PowerfulFibonacci}, while primitive-divisor theory over
polynomial rings supplies a complementary global context
\cite{CeraDaConceicao2025PrimitiveDivisors}.  Polynomial squarefreeness does
not, by itself, control the valuation of a specialized integer.  The results
below make that distinction exact.

Put $\Delta=T^2+4$.  The Binet formulas in the quadratic extension of
$\mathbb Q(T)$ give
\begin{equation}
 L_n(T)^2-\Delta F_n(T)^2=4(-1)^n,
 \qquad L_n'(T)=nF_n(T),
 \label{eq:pell-polynomial-norm-derivative}
\end{equation}
and
\begin{equation}
 \Delta F_n'(T)=nL_n(T)-TF_n(T).
 \label{eq:pell-polynomial-second-derivative}
\end{equation}
Indeed, for
$\alpha=(T+\sqrt\Delta)/2$ and $\beta=(T-\sqrt\Delta)/2$, one has
$\alpha\beta=-1$,
$F_n=(\alpha^n-\beta^n)/(\alpha-\beta)$, and
$L_n=\alpha^n+\beta^n$; subtraction and differentiation prove the three
identities.  At $T=2$ they become
\[
 L_n'(2)=nB_n,\qquad 4F_n'(2)=nA_n-B_n.
\]
Consequently, if $\ell$ is an odd prime and an odd support prime
$p\nmid\ell$ divides $B_\ell$, then $F_\ell'(2)\not\equiv0\pmod p$.
Likewise, if $q\nmid\ell$ divides $A_\ell$, then
$L_\ell'(2)\not\equiv0\pmod q$.  Coprimality of $A_\ell$ and $B_\ell$
is the only additional input.  Thus every actual repeated support factor is
a transverse specialization coincidence rather than a multiple root of the
polynomial.

\begin{lemma}[One-digit Hensel displacement]
\label{lem:pell-hensel-displacement}
Let $f\in\mathbb Z[T]$, let $p$ be prime, and suppose $p^e\mid f(t)$ with
$e\ge1$.  Then
\[
 \frac{f(t+p^eh)}{p^e}\equiv
 \frac{f(t)}{p^e}+h f'(t)\pmod p.
\]
If $p\nmid f'(t)$, there is a unique $h\pmod p$ for which
$p^{e+1}\mid f(t+p^eh)$, namely
\[
 h\equiv-\frac{f(t)}{p^e}f'(t)^{-1}\pmod p.
\]
The old representative $t$ persists precisely when
$p^{e+1}\mid f(t)$.
\end{lemma}

\begin{proof}
The integral Taylor expansion has the form
$f(t+z)=f(t)+zf'(t)+z^2G(t,z)$.  Substituting $z=p^eh$ and using
$2e\ge e+1$ gives the congruence after division by $p^e$.  Its right-hand
side is affine with nonzero slope modulo $p$, proving existence and
uniqueness of the displayed digit.  Taking $h=0$ proves the final assertion.
\end{proof}

For a simple support root write
\[
 \lambda_{f,p,e}(t)=
 -\frac{f(t)}{p^e}f'(t)^{-1}\pmod p.
\]
The product $A_\ell B_\ell$ is squarefull exactly when every level-one
displacement at the fixed parameter $T=2$ vanishes in both channels.  A
support factor of exponent at least three also has zero level-two
displacement.  This is an exact reformulation, not a proof of the required
global exclusion: the remaining problem is to prevent one fixed integer
$T=2$ from agreeing simultaneously with all of the relevant lifted roots.

Local Hensel conditions cannot supply that global prohibition when the
parameter is allowed to move.  For distinct primes $p_i$ and simple roots
$p_i\mid f_i(t_0)$, Lemma~\ref{lem:pell-hensel-displacement} gives one class
modulo $p_i^2$ for each lifted root, and the Chinese remainder theorem gives
a unique simultaneous class modulo $\prod_i p_i^2$.  At index three this
construction is completely explicit:
\[
 F_3(T)=T^2+1,\qquad L_3(T)=T^3+3T.
\]
The lifts of the $F$-channel prime $5$ and the $L$-channel prime $7$ from
$T=2$ meet at $T=282\pmod{1225}$.  At the least positive representative,
\[
 F_3(282)=79525=5^2\cdot3181,
\]
\[
 \frac{L_3(282)}2=11213307
 =3^2\cdot7^2\cdot47\cdot541,
\]
and $D=(282^2+4)/4=19882=2\cdot9941$ is squarefree.  Specializing
\eqref{eq:pell-polynomial-norm-derivative} yields the genuine global point
\[
 11213307^2-19882\,79525^2=-1.
\]
This is a full counterexample to the precise claim that simple opposite
channel roots cannot both become repeated while the parameter moves and the
Pell coefficient remains squarefree.  It has coefficient $19882$, and both
coordinates retain exponent-one factors, so it is neither a counterexample
to the fixed coefficient-two packet nor to the abc conjecture.

There is also an actual fixed-parameter collision:
\[
 F_7(T)=T^6+5T^4+6T^2+1,
\]
\[
 F_7(2)=169=13^2,
 \qquad F_7'(2)=376\equiv-1\pmod{13}.
\]
The next Hensel digit is one and
\[
 13^3\mid F_7(171)=25006385400373,
 \qquad 13^4\nmid F_7(171).
\]
Thus a simple polynomial root can specialize to an integer with exponent
two.  This refutes exactly the stronger specialization claim; it does not
give a squarefull Pell packet because $A_7=239$ occurs only to the first
power.

Finally, an exhaustive external search over positive even $T=2s$ with
$s\le10^7$ uses the unique representation $N=x^2y^3$ of a positive
squarefull integer with $y$ squarefree.  It checks $43{,}355{,}470$
canonical representations.  The sole squarefull value in the full
$F_3$ channel is
\[
 F_3(682)=465125=5^3\cdot61^2,
\]
whereas
\[
 \frac{L_3(682)}2=158608307
 =11\cdot13\cdot31\cdot37\cdot967
\]
is not squarefull.  Independent Python and C++ implementations agree.  Lean
checks the singleton identities, the nonsquarefull witness, the Pell norm,
and squarefreeness of $116282=2\cdot53\cdot1097$; the entire
$43$-million-case enumeration is not yet replayed in the kernel and is
therefore reported as an external finite computation.

The original formal companion proves the integer Hensel divisibility core,
supplied derivative readouts, index-three Taylor and CRT certificates, and the
displayed numerical points.  At that stage the all-index identities, general
Taylor and simultaneous-CRT theorems, all-support squarefull equivalence, and
one bundled moving witness were paper-level obligations.  The supplemental
Section~\ref{sec:pell-all-index-formalization} now proves those exact items in
Lean under their displayed hypotheses.  It does not supply the fixed-$T=2$
all-support transversality or the global exclusion described below.

The fixed $D=2$ route remains open.  Its smallest surviving target is to
prove that every odd prime index has a support prime with nonzero level-one
displacement at $T=2$, or more narrowly to exclude the opposite-channel
level-one/level-two packet after imposing the previously proved rank,
character, factor-quotient, all-order-tail, and endpoint-curvature
constraints.  No counterexample meeting those complete premises is known,
so difficulty does not retire the route.
